% Intended LaTeX compiler: lualatex
\documentclass[a4paper,11pt]{article}

\usepackage{amsmath}
\usepackage{fontspec}
\usepackage{graphicx}
\usepackage{longtable}
\usepackage{wrapfig}
\usepackage{rotating}
\usepackage[normalem]{ulem}
\usepackage{capt-of}
\usepackage{hyperref}
\usepackage[danish]{babel}
\usepackage[a4paper,margin=1.8cm]{geometry}
\usepackage{amsmath,array,pgffor}
\pagestyle{empty}
\renewcommand{\arraystretch}{2.2}
\newcommand{\kort}[2]{\fbox{\parbox[c][1.5cm][c]{0.45\textwidth}{\centering\small\textbf{#1}\quad #2}}}
\newcommand{\kopier}[3]{\foreach \i in {1,...,#1}{\noindent\kort{#2}{#3}\hfill\kort{#2}{#3}\par\vspace{2mm}}}
\setlength{\parindent}{0em}
\parskip 1.5ex
\date{}
\title{Facit (lærer)}
\hypersetup{
 pdfauthor={},
 pdftitle={Facit (lærer)},
 pdfkeywords={},
 pdfsubject={},
 pdfcreator={},
 pdflang={Danish}}
\begin{document}

\maketitle
\section*{Udgave 1: kort til hver linje}
\label{sec:orgb93ca6c}

\begin{center}
\begin{tabular}{|c|c|c|c|c|c|}
\hline
Linje & Ark A & Ark B & Ark C & D, p.1 & D, p.2\\
\hline
1 & K6 & K6 & K6 & K12 & K13\\
2 & K1 & K1 & K1 & K8 & K14\\
3 & K7 & K7 & K7 & K2 & K16\\
4 & K3 & K4 & K5 &  & \\
5 & K8 & K8 & K8 &  & \\
6 & K2 & K2 & K2 &  & \\
7 & K7 & K7 & K7 &  & \\
\hline
\end{tabular}
\end{center}

Overflødige kort på ark A-C: K9, K10 (falsk), K11 og K12-K16. K9 kan bruges i en alternativ version i stedet for K8 og K2, hvilket er et godt diskussionsspørgsmål.
\section*{Udgave 2: rigtig rækkefølge og kort}
\label{sec:org924371b}

Bogstaverne står i den rækkefølge, de skal stå i beviset.

\begin{center}
\begin{tabular}{|l|l|l|}
\hline
Ark & Rækkefølge (1 til 7) & Kort i den rækkefølge\\
\hline
A & D, B, G, E, A, F, C & K6, K1, K7, K3, K8, K2, K7\\
B & C, E, A, G, D, B, F & K6, K1, K7, K4, K8, K2, K7\\
C & B, F, D, A, G, C, E & K6, K1, K7, K5, K8, K2, K7\\
D, påstand 1 & B, C, A & K12, K8, K2\\
D, påstand 2 & C, A, B & K13, K14, K16\\
\hline
\end{tabular}
\end{center}
\section*{Svar til udvalgte opgaver}
\label{sec:orgf981ba6}
\begin{itemize}
\item Ark A, opg. 1: \(\log(100000)=5=2+3\).
\item Ark A, opg. 2: Er \(a=-10\), findes \(\log(a)\) ikke, så linje 1 går galt.
\item Ark A, opg. 3: \(\log(5)\approx 0{,}70\), men \(\log(2)+\log(3)\approx 0{,}78\). Potensreglen virker for produkter af potenser, ikke for summer.
\item Ark B, opg. 3: Fx \(\log(1000/10)=2\), men \(\log(1000)/\log(10)=3\).
\item Ark C, opg. 2: \(r=\frac12\), fordi \(\sqrt a=a^{1/2}\).
\item Ark C, opg. 3: Fx \(a=10\), \(r=2\): \(\log(10^2)=2\), men \((\log 10)^2=1\).
\item Ark D, opg. 2: \(\log(abc)=\log((ab)\cdot c)=\log(ab)+\log(c)=\log(a)+\log(b)+\log(c)\).
\item Ark D, opg. 3: \(\ln(a\cdot b)=\ln(a)+\ln(b)\). Grundtallet bliver \(e\), og K1, K2, K3 og K8 ændres, så \(10\) erstattes af \(e\).
\end{itemize}
\end{document}
